\section{Choosing translates directly on the ramified source}
\label{sec:affine-spanning}

The volume comparison needs a small family of translates whose images span
almost all of the target. The source first chooses translates of the entire
space of Poincar\'e series. It then restricts each map to the ramified
source and subtracts a possible dimension loss for every translate. One
can restrict at each choice instead. This gives the following stronger
version of the spanning conclusion. Here $b$ is the weight parameter,
$d$ is the target dimension, and $r$ is the ramified source dimension.
For $b\ge 8$ there is a set $J$ of translates satisfying
\begin{equation}
 |J|r\le 5d\log(1+2b^2),\qquad
 \operatorname{codim}\sum_{s\in J}
  \operatorname{im}(\text{normalized Poincar\'e map at }s)
 \le \frac{20d}{b^4}.
 \label{eq:affine-spanning-bound}
\end{equation}
The volume argument asks for codimension at most $4d/b^2$ with exactly
this translate budget. The proof below uses the existing geometric moment
estimates at their original precision.

\subsection*{An affine recurrence for the missing dimension}

Let $S$ and $V$ have positive dimensions $h$ and $d$. Let $P_s:S\to V$
be a nonempty finite family of linear maps and let $S_0\subset S$ have
dimension $r\ge (1-\tau)h$, where $0\le\tau\le1$. Suppose that for
every $F\subset V$ of codimension $c$ there is an $s$ with
\begin{equation}
 \dim P_s^{-1}(F)\le h\left(1-\frac{c}{2d}+\omega\right).
 \label{eq:affine-spread}
\end{equation}
Assume $0<h\le 2d$ and $\omega\ge0$. For each integer $u\ge 0$
one can choose at most $u$ indices whose restricted images have
codimension at most
\begin{equation}
  d\left[\kappa+(1-\kappa)\left(1-\frac{h}{2d}\right)^u\right],
  \qquad \kappa=2\omega+2\tau.
  \label{eq:affine-spread-recurrence}
\end{equation}
In particular, the missing fraction is at most
$\kappa+\exp(-uh/(2d))$.

To prove this, begin with $F=0$. If $F$ is the current span, use
\eqref{eq:affine-spread} and compose $P_s|_{S_0}$ with the quotient
$V\to V/F$. Its kernel is $S_0\cap P_s^{-1}(F)$. The dimension it adds is
therefore at least
\[
 r-\dim P_s^{-1}(F)\ \ge\ \frac{h}{2}\left(\frac{c}{d}
                                      -2\omega-2\tau\right).
\]
Writing $q_j$ for the missing fraction after $j$ choices and
$t=h/(2d)$, we have $q_{j+1}\le(1-t)q_j+t\kappa$ and $q_0=1$.
Iteration gives \eqref{eq:affine-spread-recurrence}, including when the
one-step lower bound is negative near full coverage. Finally
$(1-t)^u\le e^{-tu}$. Repeated indices contribute only one element to
the resulting set. The loss $\tau$ enters the limiting missing fraction.
It is absorbed into the fixed point of the recurrence.

\subsection*{Checking the hypotheses against the source}

Write $P$ for the scaled full-source Poincar\'e map, $e$ for orthogonal
projection onto the target, and $s$ for the unitary action of a member
of the finite group. The target is invariant and $e$ commutes with the
action. The geometric estimates in
\texttt{SpanningTheorem.lean}, at
\Source{New/Spanning/SecondMoment/SpanningTheorem}{257}{the projected near-isometry} and
\Source{New/Spanning/SecondMoment/SpanningTheorem}{288}{the normalized second moment},
give, in Hilbert--Schmidt norm,
\[
 \|eP\operatorname{pr}_U\|_{\rm HS}^2\ge\dim U-5b^{-4}h
 \quad(U\subset S),\qquad
 d\,\underset{s}{\operatorname{avg}}\,\|P^*sP\|_{\rm HS}^2
 \le(1+9b^{-4})h^2 .
\]
In particular $U=\ker(eP)$ gives $d\ge(1-5b^{-4})h$. Thus $h\le 2d$
for $b\ge8$. Write $P_s=esP$.
For a subspace $F$ of relative dimension $a=1-c/d$, put
$K=|G|^{-1}\sum_s sPP^*s^*$. Cycling the trace gives
\[
 \|K\|_{\rm HS}^2=\underset{s}{\operatorname{avg}}\,\|P^*sP\|_{\rm HS}^2,
 \qquad
 \underset{s}{\operatorname{avg}}\,\|\operatorname{pr}_F sP\|_{\rm HS}^2
 =\operatorname{tr}(\operatorname{pr}_F K)
 \le h\sqrt{a(1+9b^{-4})}.
\]
Choose $s$ whose projected energy is at most this average. On
$U=P_s^{-1}(F)$ one has $esP=\operatorname{pr}_F sP$, and unitarity
and $es=se$ give
\[
 \dim U-5b^{-4}h
 \le \|eP\operatorname{pr}_U\|_{\rm HS}^2
 \le \|\operatorname{pr}_F sP\|_{\rm HS}^2.
\]
Now $2\sqrt{a(1+9b^{-4})}\le a+1+9b^{-4}$ proves
\eqref{eq:affine-spread} with $\omega=(5+9/2)b^{-4}$.
This is the trace argument of
\Source{New/Spanning/SecondMoment/FewTranslates}{320}{the one-translate lemma}
with the two errors kept separate.

The ramified codimension estimate used in
\Source{New/Reduction/ShortAndEnough}{410}{the published choice of translates} is
$h-r\le b^{-8}h$. The source-coordinate change transports $S_0$;
the nonzero normalization factors preserve its image. Apply the
recurrence with $\tau=b^{-8}$ and
\[
        u=\left\lceil\frac{2d}{h}\log(b^5)\right\rceil .
\]
The codimension fraction is at most $19b^{-4}+2b^{-8}+b^{-5}\le20b^{-4}$. For the count, use $r\le h$, $h\le2d$ and $b\ge8$ to obtain
\[
 |J|r\le 10d\log b+h\le 10d\log b+2d
 \le 5d\log(2b^2)\le5d\log(1+2b^2).
\]
This proves \eqref{eq:affine-spanning-bound}. Since
$20/b^4\le4/b^2$ in this range, Theorem B supplies its published
interface to the determinant comparison. A minimal formal substitution can use the existing one-translate lemma verbatim: taking $\varepsilon=9b^{-4}$ gives $\omega=3\varepsilon/2$ and codimension at most $29d/b^4$, also sufficient. The $20$ uses the separate errors above. The new step is the recurrence on
the restricted image and its iteration. The stronger estimate does
not change the exponent in the prime sum, which also depends on the
length and cost of the theta chain.

\subsection*{Why the number of translates contains a logarithm}

There is a finite example with ideal moments in which the logarithm
has the correct order. Give $\mathbb C^{\binom Nk}$ coordinates indexed
by the $k$-subsets of an $N$-element set. The symmetric group permutes
the coordinates. Let $P$ include the coordinate subspace spanned by
the $k$-subsets containing a specified point. Here $d=\binom Nk$,
$h=\binom{N-1}{k-1}$, and $h/d=k/N$. The translated image picks a
point and spans all coordinates containing it. Each coordinate occurs
in $k$ of the $N$ images, so the average of their orthogonal projectors
is $(h/d)I$. In particular the first moment is exact and the mean of
$\|P^*sP\|_{\rm HS}^2$ is $h^2/d$.

Choosing $j$ distinct points leaves exactly
$\binom{N-j}{k}/\binom Nk$ of the coordinates uncovered. For fixed $L>0$,
take $k\to\infty$, $k/N\to0$, and $j=\lfloor LN/k\rfloor$. The logarithm
of this fraction is
\[
 \sum_{a=0}^{k-1}\log\left(1-\frac{j}{N-a}\right)=-L+o(1).
\]
Indeed the sum of the fractions tends to $L$ and the sum of their
squares tends to zero. Thus even exact moments can require order
$(d/h)\log(1/\theta)$ translates to reach a prescribed missing
fraction $\theta$. This example concerns the information in the two
moment bounds. The particular Poincar\'e maps could have additional
structure.
